% =====================================================================
% chapters/01_introduction_and_pm.tex
% Content only -- \chapter{Process Mining} is issued in main.tex.
% =====================================================================
Digital transformation is reshaping organizational functions across virtually every domain and industry~\cite{verhoef2021digitaltransformation}. As a consequence, nearly every operational step leaves a footprint that can be traced: a record of when, where and by whom each activity was executed.
Process mining builds exactly on this phenomenon, reconstructing the real world execution of the processes from the digital traces they generate.~\cite{aalst2011manifesto} \\
This chapter introduces the foundational concepts to formalize this reconstruction objective. 

\section{Types of Process Mining} 
In its embrional form, three types of process mining are theoreticized. 
	\begin{itemize}[leftmargin=*]
		\item \textbf{Discovery}: taking an event log and producing a model without using a-priori information ;
		\item \textbf{Conformance}:an existing process model is compared with the one reconstructed with event logs to check if the empirical records conforms to the model and vice versa;
		\item \textbf{Enhancement}: the idea is to improve an already existing process model using information about the actual process, as recorded in some event log.
	\end{itemize}
	In more recent works, the Enhancement type of process mining is further developed in process \emph{Reengineering} and \emph{Operational support}.
	



\section{Definition of an Event Log} 
A business process consists of a sequence of activities performed to achieve a specific outcome. An \emph{event log} is defined as a collection of events recording the
execution of such a process. Each event refers to a single activity, performed at a specific point in time, and is associated with a unique case~\cite{reinkemeyer2020nutshell}. 
In its minimal form, an event log requires the following three attributes for every event:
	\begin{itemize}[leftmargin=*]
		\item \textbf{Case ID}: a unique identifier for the case to which the event belongs;
		\item \textbf{Activity}: a label specifying which activity was
		executed;
		\item \textbf{Timestamp}: the time at which the activity was
		performed.
	\end{itemize}

Additional information about events might be stored in the event logs. In fact, modern process mining techniques incorporate different types of extra data elements in the process reconstruction~\cite{aalst2012opportunities}.

\subsection{Formalizing the Event Log}
\label{subsec:event-log-formal}

An event log can be understood as a particular view on a broader
\emph{event collection}. While an event collection has no fixed notion of
case, an event log groups events into cases, giving each one a well-defined
structure. We adopt the formalization of~\cite{aalst2018reconstruction}.

\begin{definition}[Event Collection]
	\label{def:event-collection}
	Let $\mathcal{E}$ be the universe of event identifiers, $\mathcal{U}_{Att}$
	the universe of attribute names, and $\mathcal{U}_{Val}$ the universe of
	attribute values. An \emph{event collection} is a triple
	$EC = (E, Att, \pi)$ with $E \subseteq \mathcal{E}$,
	$Att \subseteq \mathcal{U}_{Att}$, and
	$\pi \in E \to (Att \not\to \mathcal{U}_{Val})$. For an event $e \in E$ and
	an attribute $a \in \mathrm{dom}(\pi(e))$, $\pi(e)(a)$ denotes the value of
	attribute $a$ for event $e$.
\end{definition}

Extracting an event log from an event collection requires selecting a
subset of events and mapping their attributes onto the three mandatory
process mining concepts: case, activity, and timestamp. Every trace in the
resulting log is additionally assumed to begin with a unique artificial
start activity $\triangleright$ and end with a unique artificial end
activity $\square$, which simplifies the formal treatment of process
discovery without any loss of generality.

\begin{definition}[Event Log]
	\label{def:event-log}
	Let $A$ be a set of activities with $\{\triangleright, \square\} \subseteq A$,
	and let $\mathcal{B}(A^{*})$ denote the set of all finite multisets of
	sequences over $A$. An \emph{event log} is a non-empty multiset of traces
	$L \in \mathcal{B}(A^{*})$. A trace $\sigma \in L$ is a sequence of
	activities of the form
	\[
	\sigma = \langle \triangleright, a_1, a_2, \dots, a_n, \square \rangle,
	\qquad \{\triangleright, \square\} \cap \{a_1, \dots, a_n\} = \emptyset.
	\]
	The set of activities occurring in $L$ is denoted
	$\mathrm{act}(L) = \bigcup_{\sigma \in L} \{a \in \sigma\} \subseteq A$.
\end{definition}

As a concrete illustration, the log
$L_1 = [\langle \triangleright, a, b, c, d, \square \rangle^{45},
\langle \triangleright, a, c, b, d, \square \rangle^{35},
\langle \triangleright, a, e, d, \square \rangle^{20}]$
contains $100$ traces and $580$ events in
total~\cite{aalst2018reconstruction}, with
$\mathrm{act}(L_1) = \{\triangleright, a, b, c, d, e, \square\}$.

\subsection{Log and Model Abstractions}
\label{subsec:abstractions}

Process discovery techniques rarely operate directly on an event log.
Instead, most approaches first reduce the log to a more compact
\emph{abstraction}, discover a corresponding abstraction on the model side,
and only then reconstruct a process model. This two-sided mapping is what
allows discovery algorithms with very different internal representations
(directly-follows graphs, transition systems, declarative constraints) to be
described within a single, unifying framework~\cite{aalst2018reconstruction}.

\begin{definition}[Log and Model Abstraction Functions]
	\label{def:abstractions}
	Let $L \in \mathcal{B}(A^{*})$ be an event log and
	$PM \in \mathcal{P}(A^{*})$ a process model, i.e., a set of traces over
	activity set $A$. We distinguish three functions:
	\begin{itemize}[leftmargin=*]
		\item $m_{L2LA}$, the \emph{log abstraction function}, maps an event log
		onto a \emph{log abstraction} $m_{L2LA}(L)$ (e.g., a weighted
		directly-follows graph);
		\item $m_{LA2MA}$, applied to a log abstraction, yields a
		\emph{model abstraction} $m_{LA2MA}(m_{L2LA}(L))$, typically by
		discarding infrequent behavior through a threshold;
		\item $m_{PM2MA}$, the \emph{model abstraction function}, maps a process
		model onto the same type of abstraction, so that $m_{PM2MA}(PM)$ is
		directly comparable with $m_{LA2MA}(m_{L2LA}(L))$.
	\end{itemize}
\end{definition}

Recovering a process model from a model abstraction requires inverting
$m_{PM2MA}$. Since multiple models may share the same abstraction, this
inverse is resolved by assuming $m_{PM2MA}^{-1}$ returns a canonical model
for a given abstraction, which yields the overall discovery expression
\begin{equation}
	\label{eq:discovery-composition}
	PM = m_{PM2MA}^{-1}\bigl(m_{LA2MA}(m_{L2LA}(L))\bigr),
\end{equation}
turning an event log directly into a process model~\cite{aalst2018reconstruction}.
Equation~\eqref{eq:discovery-composition} will serve as the organizing
principle for Chapter~\ref{chap:graphtheory}: the directly-follows graph
introduced in Section~\ref{sec:dfg} is precisely one instance of the log
abstraction $m_{L2LA}$, and the arborescence, shortest-path, and min-cut
algorithms discussed there each operate downstream of this abstraction step.
	
\subsection{The Directly-Follows Graph}
\label{subsec:dfg-formal}

A \emph{directly-follows graph} (DFG) is an \emph{abstraction} of the
behavior actually observed in an event log: it records only whether one
activity is, with what frequency, immediately followed by another, without
any of the causal or concurrency semantics a formal process model such as a
Petri net would provide. The DFG is accordingly the simplest and most
widely used log abstraction $m_{L2LA}$ introduced in
Definition~\ref{def:abstractions}, and is the basis of nearly every
commercial process mining tool~\cite{aalst2018reconstruction}.

\begin{definition}[Directly-Follows Graph]
	\label{def:dfg}
	Let $A$ be a set of activities with $\{\triangleright, \square\} \subseteq A$.
	A \emph{directly-follows graph} is a pair $DFG = (A, R)$ with
	$R \subseteq A \times A$ a relation such that no activity directly follows
	$\square$ and no activity directly precedes $\triangleright$, i.e.,
	\[
	\{ (a, \triangleright) \in R \mid a \in A \} =
	\{ (\square, a) \in R \mid a \in A \} = \emptyset.
	\]
\end{definition}

In practice, what matters is not merely whether one activity directly
follows another, but \emph{how often} it does so; this motivates a weighted
variant obtained directly from an event log via the log abstraction function
$m_{L2LA}$ of Definition~\ref{def:abstractions}.

\begin{definition}[Weighted Directly-Follows Graph]
	\label{def:wdfg}
	Let $L \in \mathcal{B}(A^{*})$ be an event log. The \emph{weighted
		directly-follows graph} of $L$ is $WDFG = (A, R_w)$, where $A = \mathrm{act}(L)$
	and
	\[
	R_w = \bigcup_{\sigma \in L} \bigl[\, (a, b) \in A \times A \mid
	\exists_{1 \leq i < |\sigma|}\ \sigma(i) = a \ \wedge\ \sigma(i+1) = b \,\bigr].
	\]
	For $(a, b) \in A \times A$, $R_w(a, b)$ denotes the number of times activity
	$a$ is directly followed by activity $b$ across all traces of
	$L$~\cite{aalst2018reconstruction}.
\end{definition}

Computing $WDFG$ from $L$ corresponds exactly to applying $m_{L2LA}$
(Definition~\ref{def:abstractions}) with the directly-follows graph as the
chosen abstraction. An ordinary, non-weighted $DFG$ is then recovered by
applying $m_{LA2MA}$: an arc $(a, b)$ is retained if and only if its weight
$R_w(a, b)$ exceeds some threshold $\theta$, absolute or relative to the
other weights in $R_w$. Crucially, Definition~\ref{def:dfg} imposes
\emph{no acyclicity constraint} on $R$: a $WDFG$ is a general weighted
digraph that may contain self-loops and longer cycles, reflecting rework and
repeated behavior in the underlying process. This property is the starting
point for the discussion in Section~\ref{sec:dfg-to-digraph}, where the
$WDFG$ is treated as a directed graph $G = (V, E, w)$ and subjected to the
graph-theoretic algorithms developed in Chapter~\ref{chap:graphtheory}.

\subsection{Petri Nets}
\label{subsec:petri-nets}

The $WDFG$ of Definition~\ref{def:wdfg} abstracts \emph{observed} behavior,
but is not itself a process model: it makes no claim about which traces are
admissible, only about which activities were seen to follow one another. A
\emph{Petri net} is the formal target model into which a (filtered) DFG is
ultimately converted by a discovery algorithm, and is accordingly the
notation used throughout this work to represent
$PM \in \mathcal{P}(A^{*})$.

\begin{definition}[Labeled Petri Net]
	\label{def:petri-net}
	A \emph{Petri net} is a tuple $N = (P, T, F)$, where $P$ is a finite set of
	places, $T$ is a finite set of transitions, $P \cap T = \emptyset$, and
	$F \subseteq (P \times T) \cup (T \times P)$ is the flow relation. A
	\emph{labeled} Petri net $N = (P, T, F, l, A)$ extends $N$ with an activity
	set $A$ and a labeling function $l \in T \not\to A$, where a transition
	without a label is \emph{silent}. Equipped with an initial and a final
	marking, the behavior of $N$ is the set of activity sequences obtainable
	by firing transitions from the initial to the final marking, i.e., a
	process model $PM \subseteq A^{*}$ in the sense of
	Definition~\ref{def:abstractions}~\cite{aalst2018reconstruction}.
\end{definition}

\subsection{Process discovery: from Event Log to Process Model}
\label{subsec:dfg-to-petri}

The \texttt{pm4py} pipeline implemented in \texttt{loader.py} follows
precisely the composition of Equation~\eqref{eq:discovery-composition}.
Given an event log $L$, \texttt{pm4py.discover\_dfg} computes the weighted
directly-follows graph $m_{L2LA}(L) = WDFG$ of Definition~\ref{def:wdfg}. A
discovery algorithm then applies $m_{LA2MA}$, typically by thresholding
$R_w$ to discard infrequent arcs, before applying $m_{PM2MA}^{-1}$ to
synthesize a labeled Petri net $N$ (Definition~\ref{def:petri-net}) whose
behavior $PM_N$ reproduces the retained directly-follows relations. The
resulting net is the process model used in
Chapter~\ref{chap:comparativeanalysis} to compare the empirical structure of
the Road Traffic Fine Management and Sepsis Cases event
logs~\cite{aalst2018reconstruction}.

\subsection{The Alpha Algorithm}
\label{subsec:alpha-algorithm}

The $m_{PM2MA}^{-1}$ step of Equation~\eqref{eq:discovery-composition} is,
in practice, carried out by a concrete discovery algorithm. The
\emph{alpha algorithm} ($\alpha$-algorithm) was the first discovery
technique to take concurrency as an explicit starting point, rather than as
an afterthought layered on top of a sequential
model~\cite{aalst2004workflow, aalst2018reconstruction}. Given an event log
$L$, the $\alpha$-algorithm derives, for every pair of activities
$a, b \in \mathrm{act}(L)$, one of four ordering relations from the
directly-follows information of Definition~\ref{def:wdfg}:

\begin{itemize}[leftmargin=*]
	\item $a \to_L b$ (\emph{causality}): $a$ is directly followed by $b$
	somewhere in $L$, but $b$ is never directly followed by $a$;
	\item $a \parallel_L b$ (\emph{parallelism}): $a$ is directly followed by
	$b$ \emph{and} $b$ is directly followed by $a$ somewhere in $L$, indicating
	concurrent execution;
	\item $a \#_L b$ (\emph{independence/choice}): neither $a$ directly
	follows $b$ nor $b$ directly follows $a$ anywhere in $L$;
	\item $a \leftarrow_L b$: the symmetric counterpart of $a \to_L b$.
\end{itemize}

These relations are used to group activities into sets that become the
input and output places of a labeled Petri net
(Definition~\ref{def:petri-net}): causal relations $a \to_L b$ are turned
into places connecting $a$ to $b$, while parallel and choice relations
determine how multiple outgoing or incoming places are merged. The
resulting net is, by construction, the canonical model that
$m_{PM2MA}^{-1}$ returns for the (filtered) directly-follows abstraction
$m_{LA2MA}(m_{L2LA}(L))$ of a log $L$~\cite{aalst2004workflow}.

The $\alpha$-algorithm is guaranteed to rediscover a behaviorally
equivalent process model for a restricted, well-defined class of
processes~\cite{aalst2004workflow}. Outside this class, in particular in
the presence of incomplete logs, short loops, or non-free-choice
constructs, the algorithm may return a model that is not
sound (e.g., exhibiting deadlocks) or that fails to allow for behavior
actually observed in $L$~\cite{aalst2018reconstruction}. These limitations
motivate the family of more robust discovery techniques, discussed next,
that \texttt{pm4py} implements alongside the
$\alpha$-algorithm~\cite{berti2019process}.

\section{Spaghetti Processes, Quantified}
\begin{itemize}[leftmargin=*]
	\item \textbf{Core Concept:} Reframe ``spaghetti'' from a descriptive term into a measurable graph property.
	\item \textbf{Candidate Metrics} (choose one as the working definition):
	\begin{itemize}
		\item Edge density relative to $|V|$
		\item Mean out-degree
		\item Sparsity index
	\end{itemize}
	\item \textbf{Dataset Application:} To be computed for both Chapter 4 datasets (\emph{Road Traffic Fine Management}, \emph{Sepsis Cases})—gives a concrete before/after number once Edmonds' cleaning (Ch.~2) is applied.
\end{itemize}

\section{Cycles as a Natural Part of Process Flow}
\begin{itemize}[leftmargin=*]
	\item \textbf{Self-loop:} Edge $(v,v)$.
	\item \textbf{Rework loop:} Any cycle of length $\geq 2$ in the DFG.
	\item \textbf{Rework Ratio Statistic:}
	\[
	\text{Rework Ratio} = \frac{\text{loop-edge weight}}{\text{total edge weight}}
	\]
	(another figure to carry into Chapter 4).
	\item \textbf{Open Question (carried into Ch.~2):} Does any downstream algorithm require this structure to be acyclic, and if so, at what cost?
\end{itemize}